\documentclass[11pt]{article}
\usepackage{fontspec}
\usepackage[french]{babel}
\usepackage{amsmath}
\usepackage{amssymb}
\usepackage{amsthm}
\usepackage{mathrsfs}
\usepackage{tikz-cd}
\usepackage{enumitem}
\usetikzlibrary{babel}

\newcommand{\printedpage}[1]{\marginpar{\footnotesize p.~#1}}
\newcommand{\msp}[1]{\marginpar{\footnotesize [#1]}}
\newcommand{\uHom}{\underline{\mathrm{Hom}}}

\begin{document}

\printedpage{43}
% head: A. Grothendieck : Dérivateurs, Ch. XII, Édition des Universités de Montpellier II et Paris VII ; foot: Version 7 décembre 2002

\msp{page 59}

\section*{V\quad Sommes amalgamées et carrés $W$-cocartésiens dans $\mathrm{Cat}$}

\subsection*{1\quad Généralités sur $U \overset{i}{\hookrightarrow} X \overset{j}{\hookleftarrow} T$, et les sommes amalgamées.}

\subsubsection*{1.1\quad Données de recollement entre un ouvert et un fermé.}

Soit $X$ dans Cat, $U$ un ouvert dans $X$, $T$ le complémentaire fermé. (\textbf{NB} Les ouverts correspondent bien aux ouverts de $\mathrm{Ob}\,X$, topologisé par sa relation de préordre naturelle $x \leq y \iff \mathrm{Hom}_X(x,y) \neq \emptyset$, de même pour les fermés. Les topologies associées aux relations de préordre sont celles où toute intersection d'ouverts est encore ouverte, i.e. toute réunion de fermés est encore fermée. Donc les ouverts d'une telle topologie sont les fermés d'une topologie duale, correspondant à la relation de préordre opposée. La relation de préordre associée à la topologie est
\[
x \leq y \qquad\qquad \iff \qquad\qquad x \in \overline{\{y\}},
\]
la topologie associée au préordre sur $I$ est celle pour laquelle les ouverts sont les réunions d'ensembles de la forme $I/x$ (lesquels forment une base d'ouverts), ou encore les \textit{`cribles'} ($x \in U$, $y \leq x \Longrightarrow y \in U$), les fermés sont les réunions d'ensembles $x\backslash I$, c'est-à-dire les cocribles. Pour toute partie $A$ de $I$, il y a un plus petit ouvert contenant $A$, qui est aussi son plus petit voisinage, c'est la réunion des $I/x$ pour $x \in A$. Et dualement il y a un plus grand fermé contenu dans $A$, dont le complémentaire est le plus petit ouvert contenant $B = I\backslash A$, donc $\bigcup_{x\in B} I/x$. L'adhérence de $A$ est la réunion des adhérences $x\backslash I$ de ses points.)

\msp{page 60}

La situation $(X,U,T)$ se reconstitue, à isomorphisme unique près, par la connaissance de

\begin{enumerate}[label={}]
\item[$1^\circ$)] Les deux catégories $U$, $T$.
\item[$2^\circ$)] Le bifoncteur $\displaystyle\left\{\begin{aligned}
&H : U^\circ \times T \longrightarrow \mathrm{Ens}\\[1ex]
&H(u,t) = \mathrm{Hom}_X(u,t).
\end{aligned}\right.$
\end{enumerate}

De plus, ces données $1^\circ$) $2^\circ$) en sont soumises à aucune restriction. On a, plus exactement % [sic] en sont -> ne sont, negation required by the sense

\textbf{Proposition 1.} \textit{Un isomorphisme} \texttt{[plutôt une équivalence de catégories]} \textit{entre la catégorie $\mathcal{T}$ dont les objets sont les triples $(X,U,T)$, avec $X \in \mathrm{Ob}\,\mathrm{Cat}$, $U \in \mathrm{Ouv}(X)$, $T \in \mathrm{Ferm}(X)$, $U$ et $T$ complémentaires l'un de l'autre, les morphismes de $(X,U,T)$ dans $(X',U',T')$ étant les foncteurs $f : X \longrightarrow X'$ qui induisent $U \longrightarrow U'$, $T \longrightarrow T'$ (d'où
\printedpage{44}
% head: A. Grothendieck : Dérivateurs, Ch. XII, Édition des Universités de Montpellier II et Paris VII ; foot: Version 7 décembre 2002
$U = f^{-1}(U')$, $T = f^{-1}(T')$), avec la catégorie $\mathcal{T}'$ des triples $(U,T,H)$, de deux catégories $U$, $T$ et d'un bifoncteur}
\[
H : U^\circ \times T \longrightarrow \mathrm{Ens}
\]
\hspace*{2em}\textit{(ou encore un préfaiscau $U^\circ \longrightarrow T^\vee$ sur $U$, donc un objet de} $\uHom(U^\circ,T^\vee)$\textit{)),} % [sic] préfaiscau -> préfaisceau, misspelling

\textit{un morphisme de $(U,T,H)$ dans $(U',T',H')$ étant un triple}
\[
(f_U, f_T, u) \qquad\qquad
\left\{\begin{array}{l}
\left.\begin{array}{llll}
f_U : & U & \longrightarrow & U' \\[1ex]
f_T : & T & \longrightarrow & T'
\end{array}\right\}\text{des foncteurs}\\[2ex]
u : H \longrightarrow H' \circ (f_U^\circ \times f_T),\\
\text{i.e. } H(u,t) \longrightarrow H'(f_U(u), f_T(t))\\
\text{un homomorphisme fonctoriel}\\
\text{(flèche dans } (U \times T^\circ)^\wedge \simeq \uHom(U^\circ,T^\vee)\text{)),}
\end{array}\right.
\]
\textit{composition des flèches évidente.}

Si $U \hookrightarrow X \hookleftarrow T$ est un triple, et
\[
f_{U'} : U' \longrightarrow U,\ f_{T'} : T' \longrightarrow T
\]
des flèches dans Cat, il y a une façon privilégiée de `recoller' $U'$ et $T'$, de façon à ce que l'objet $(X',U',T')$ s'envoie dans $(X,U,T)$ en induisant

\msp{page 61}

$f_{U'}$, $f_{T'}$:
\[
\begin{tikzcd}[column sep=large, row sep=large]
U' \arrow[r, hook, "i'"] \arrow[d, "f_{U'}"'] & X' \arrow[d, "f"] & T' \arrow[l, hook', "j'"'] \arrow[d, "f_{T'}"] \\
U \arrow[r, hook, "i"] & X & T, \arrow[l, hook', "j"']
\end{tikzcd}
\]
elle s'obtient à l'aide des bifoncteurs
\[
\left\{\begin{aligned}
&H'(u',t') \overset{\text{déf}}{=} H(f_{U'}(u'), f_{T'}(t'))\\
&\text{i.e. } H' = H \circ (f_{U'}^\circ \times f_{T'}).
\end{aligned}\right.
\]
Ce triple, i.e. cet objet de $\mathcal{T}$, représente le foncteur suivant sur $\mathcal{T}$ :
\[
(X'',U'',T'') \longmapsto \mathrm{Hom}_{\mathrm{Cat}}(U'',U) \times \mathrm{Hom}_{\mathrm{Cat}}(T'',T),
\]
i.e. on a pour tout objet $(X'',U'',T'')$ de $\mathcal{T}$ une bijection canonique, fonctorielle en cet objet
\[
\mathrm{Hom}_{\mathcal{T}}((X'',U'',T''),(X',U',T')) \simeq \mathrm{Hom}_{\mathrm{Cat}}(U'',U) \times \mathrm{Hom}_{\mathrm{Cat}}(T'',T). \tag*{$(*)$}
\]

\printedpage{45}
% head: A. Grothendieck : Dérivateurs, Ch. XII, Édition des Universités de Montpellier II et Paris VII ; foot: Version 7 décembre 2002

Nul doute qu'on a même
\[
\uHom_{\mathcal{T}}((X'',U'',T''),(X',U',T')) \simeq \uHom_{\mathrm{Cat}}(U'',U) \times \uHom_{\mathrm{Cat}}(T'',T), \tag*{$(**)$}
\]
tout aussi fonctoriel, dont $(*)$ se déduit en passant aux ensembles d'objets. On dit que le triple $(X',U',T')$ est le triple \textit{induit} par le triple $(X,U,T)$, via les foncteurs $f_{U'}$, $f_{T'}$.

\textbf{NB} Si $f_{U'}$, $f_{T'}$ sont pleinement fidèles, alors il en est de même de $f : X' \longrightarrow X$. Nous utiliserons cette opération de recollement privilégié surtout quand $U' = U$, $f_{U'} = \mathrm{id}_U$, ou inversement $T' = T$, $f_{T'} = \mathrm{id}_T$.

\subsubsection*{1.2\quad Relation avec le recollement des topos.}

La situation d'un triple
\[
(\mathcal{X},\mathcal{U},\mathcal{T})
\]
d'un topos $\mathcal{X}$, d'un ouvert $\mathcal{U}$ de $\mathcal{X}$, et d'un

\msp{page 62}

sous-topos fermé complémentaire, est toute similaire (\textbf{NB} Si $\mathcal{X} = \mathrm{Top}(X)$, avec $X$ dans Cat, alors $\mathrm{Ouv}(\mathcal{X}) \simeq \mathrm{Ouv}(X)$, puisque les ouverts correspondent aux sous-faisceaux du faisceau final sur $\mathcal{X}$, i.e. du préfaisceau final sur $X$. Cette bijection est compatible avec le passage aux complémentaires.) La théorie des topos dit que les couples $(\mathcal{X},\mathcal{U},\mathcal{T})$ \texttt{[plutôt triples]} se reconstituent, à équivalence définie à isomorphisme unique près, à l'aide d'un triple
\[
\left\{\begin{aligned}
&\mathcal{U},\ \mathcal{T} \text{ deux topos}\\[1ex]
&\varphi : \mathrm{Fais}\,(\mathcal{U}) \longrightarrow \mathcal{T} \text{ un foncteur exact à gauche (`accessible'),}
\end{aligned}\right.
\]
ce foncteur $\varphi$ étant obtenu, quand on part de
\[
\mathcal{U} \xrightarrow{\;i\;} \mathcal{X} \xleftarrow{\;j\;} \mathcal{T},
\]
comme
\[
\varphi = j^*i_*.
\]
De plus, le triple $(\mathcal{U},\mathcal{T},\varphi)$ peut être choisi arbitrairement, avec la seule restriction (pratiquement anodine, car toujours satisfaite pour tous les foncteurs naturels qu'on rencontre) que $\varphi$ soit (en plus d'être exact à gauche) `accessible'.

Dans le cas actuel $\mathcal{X} = \mathrm{Top}(X)$, $\mathcal{U} = \mathrm{Top}(U)$, $\mathcal{T} = \mathrm{Top}(T)$, le foncteur $\varphi = j^*i_*$ est non seulement exact à gauche, mais même il commute aux $\varprojlim$ quelconques, puisqu'il en est ainsi (toujours) de $i_*$, et aussi (parce que le $j$ topossique est induit par le $j : T \longrightarrow X$ catégorique) de $j^*$. En fait, on
\printedpage{46}
% head: A. Grothendieck : Dérivateurs, Ch. XII, Édition des Universités de Montpellier II et Paris VII ; foot: Version 7 décembre 2002
\msp{page 63}
trouve ceci :

\textbf{Proposition 2 :} \textit{Considérons une situation de recollement de deux topos $\mathcal{U}$, $\mathcal{T}$ par $\varphi : \mathrm{Fais}\,\mathcal{U} \longrightarrow \mathrm{Fais}\,\mathcal{T}$, foncteur exact à gauche. Supposons que $\mathcal{U} = \mathrm{Top}(U)$, $\mathcal{T} = \mathrm{Top}(T)$, avec $U$ et $T$ dans $\mathrm{Cat}$. Pour que le topos recollé $\mathcal{X}$ soit équivalent à un topos $\mathrm{Top}(X)$, avec $X$ dans $\mathrm{Cat}$, il faut et il suffit que le foncteur de recollement $\varphi$ commute aux $\varprojlim$ quelconques, ou ce qui revient au même, qu'il ait un adjoint à gauche :}
\[
\psi : \underbrace{\mathrm{Fais}\,\mathcal{T}}_{= T^\wedge} \longrightarrow \underbrace{\mathrm{Fais}\,\mathcal{U}}_{= U^\wedge}.
\]

On s'attend dès lors que la donnée de recollement $\varphi$, ou ce qui revient au même le foncteur $\psi$ ayant un adjoint à droite, i.e. commutant aux $\varinjlim$ quelconques, revient à celle d'une donnée de recollement $H$
\[
H : U^\circ \times T \longrightarrow \mathrm{Ens}
\]
pour $U$ et $T$. Mais on sait bien que les foncteurs $\psi : T^\wedge \longrightarrow \mathcal{C}$ dans une catégorie $\mathcal{C}$ avec $\varinjlim$ quelconques, qui commutent auxdites $\varinjlim$, forment une catégorie équivalente (par le foncteur $\psi \longmapsto \psi|T(\subseteq T^\wedge)$) à celle des foncteurs
\[
\psi_0 : T \longrightarrow \mathcal{C} \quad\ (= \underbrace{U^\wedge}_{\uHom(U^\circ,\mathrm{Ens})} \text{ en l'occurrence}),
\]
\msp{page 64}
et la donnée d'un tel foncteur $\psi_0$ revient bien à celle d'un $H$. Donc

\textbf{Corollaire.} \textit{La catégorie des données de recollement topossique $\varphi : U^\wedge \longrightarrow T^\wedge$ entre} $\mathrm{Top}(U)$, $\mathrm{Top}(T)$\textit{, telles que le topos recollé soit un `topos} paradisiaque\textit{', i.e. équivalent à un topos de la forme $X^\wedge$ (avec $X$ dans} Cat\textit{), est équivalente à celle des données de recollement $H : U^\circ \times T \longrightarrow$} Ens \textit{entre $U$ et $T$. La donnée topossique $\varphi$ associée à $H$ est le foncteur adjoint à droite de $T^\wedge \longrightarrow U^\wedge$ déduit de $T \xrightarrow{H} U^\wedge$ par prolongement commutant aux $\varinjlim$, ou encore elle est donnée par}
\[
\varphi(\underbrace{F}_{\in U^\wedge})(t) \simeq \mathrm{Hom}_{U^\wedge}(H(t),F),
\]
\textit{où dans le second membre $H$ est interprété comme un foncteur $T \longrightarrow U^\wedge$. Le topos recollé par $\varphi$ est canoniquement équivalent à} $\mathrm{Top}(X)$\textit{, où $X$ est l'objet de} Cat \textit{obtenu par recollement à l'aide de $H$.}

Du point de vue de la théorie des faisceaux, cela donne l'interprétation de $H$, comme exprimant le foncteur
\[
\varphi = j^*i_* : U^\wedge \longrightarrow T^\wedge
\]
par la formule précédente.

\printedpage{47}
% head: A. Grothendieck : Dérivateurs, Ch. XII, Édition des Universités de Montpellier II et Paris VII ; foot: Version 7 décembre 2002

\msp{page 65}

\subsubsection*{1.3\quad Flèches dans $\mathrm{Cat}$, de source ou but un objet recollé $(X,U,T)$.}

Quand $(X,U,T)$ est donné par une donnée de recollement
\[
U,\ T,\ H : U^\circ \times T \longrightarrow \mathrm{Ens},
\]
les flèches
\[
f : X' \longrightarrow X
\]
dans Cat se décrivent immédiatement ainsi : la donnée de $f$ définit d'abord
\[
U' = f^{-1}(U), \qquad\qquad T' = f^{-1}(T),
\]
couple d'un ouvert de $X'$ et du fermé complémentaire, de sorte que cette donnée équivaut à celle de

\begin{enumerate}[label={}]
\item[$1^\circ$)] un couple $(U',T')$ dans $X'$
\item[$2^\circ$)] un morphisme (dans $\mathcal{T}$) de $(X',U',T')$ dans $(X,U,T)$, ce qui équivaut par 1.1 à un triple
\[
\left\{\begin{array}{rccl}
f_{U'} : & U' & \longrightarrow & U\\[1ex]
f_{T'} : & T' & \longrightarrow & T\\[1ex]
u : & H_{U',T'} & \longrightarrow & H \circ (f_{U'}^\circ \times f_{T'}),
\end{array}\right.
\]
où on a posé
\[
\left\{\begin{aligned}
&H_{U',T'} : U^{\prime\circ} \times T' \longrightarrow \mathrm{Ens}\\[1ex]
&H_{U',T'}(u',t') = \mathrm{Hom}_{X'}(u',t'),
\end{aligned}\right.
\]
i.e. $H_{U',T'}$ est la donnée de recollement qui conrrespond au triple $(X',U',T')$ dans $\mathcal{T}$. % [sic] conrrespond -> correspond, misspelling
\end{enumerate}

\msp{page 66}

Nous serons plus intéressés par la situation inverse, i.e. par la question de déterminer les foncteurs
\[
\underbrace{X}_{= X_{U,T,H}} \quad \longrightarrow X'
\]
de $X$, obtenu par $(U,T,H)$, dans un objet quelconque de Cat. On trouve tautologiquement

\textbf{Proposition 3.} $\mathrm{Hom}_{\mathrm{Cat}}(X,X')$ \textit{est en correspondance bijective canonique, fonctorielle en $X'$ dans} Cat\textit{, avec l'ensemble des triples}

\end{document}
